\documentclass[aspectratio=169]{beamer}
\usetheme{Madrid}
\usecolortheme{beaver}
\usepackage[utf8]{inputenc}
\usepackage[T1]{fontenc}
\usepackage{tikz}
\usepackage{amssymb}
\usepackage{graphicx}
\setbeamertemplate{navigation symbols}{}
\setbeamercolor{title}{fg=white,bg=cyan!60!black}
\setbeamercolor{frametitle}{fg=white,bg=cyan!50!black}
\setbeamercolor{block title}{bg=cyan!15,fg=black}
\setbeamercolor{block body}{bg=cyan!5,fg=black}

\title{G7 — Game States}
\subtitle{Lesson 2: Switching States}
\author{ByteForge --- Grade 7 --- Game Maker}
\date{}

\begin{document}

\begin{frame}
    \titlepage
\end{frame}

\begin{frame}{A condition that changes the state itself}
    \Large
    The game changes its own state variable based on a condition — this self-modifying behaviour is exactly how real games transition between screens.
\end{frame}

\begin{frame}{Your Turn}
    \Large
    Open the student map, find \textbf{G7 — Game States}, and complete
    the \textbf{Switching States} challenge --- everything you just saw, they now do themselves.
\end{frame}

\end{document}
